%% kapa: exam=1 round=28 year=2017 period=1 form=A subject=real_estate subject_ko=부동산학원론 seq=25 number=105
%% source: Q-Net 감정평가사 기출문제 (cst00309) artlSeq=5199603 file=extracted/1008/5199603/2017년 제28회 감정평가사 1차시험 문제지/1교시 A형.pdf
%% conversion: kordoc 4.24.51 -> tools/gen1.py (auto); review=proofread-sonnet; flags=none
\begin{question}{105}
감정평가에 관한 규칙의 내용으로 옳지 않은 것은?
\begin{choices}
  \item 대상물건에 대한 감정평가액은 시장가치를 기준으로 결정하나, 감정평가 의뢰인이 요청하는 경우 등에는 시장가치 외의 가치를 기준으로 결정할 수 있다.
  \item 적정한 실거래가는 부동산 거래신고에 관한 법률에 따라 신고된 실제 거래가격으로서 거래 시점이 도시지역은 3년 이내, 그 밖의 지역은 5년 이내인 거래가격 중에서 감정평가업자가 인근지역의 지가수준 등을 고려하여 감정평가의 기준으로 적용하기에 적정하다고 판단하는 거래가격을 말한다.
  \item 가치형성요인은 대상물건의 경제적 가치에 영향을 미치는 일반요인, 지역요인 및 개별요인 등을 말한다.
  \item 시장가치는 감정평가의 대상이 되는 토지등(이하 "대상물건")이 통상적인 시장에서 충분한 기간 동안 거래를 위하여 공개된 후 그 대상물건의 내용에 정통한 당사자 사이에 신중하고 자발적인 거래가 있을 경우 성립될 가능성이 가장 높다고 인정되는 대상물건의 가액을 말한다.
  \item 유사지역은 감정평가의 대상이 된 부동산이 속한 지역으로서 부동산의 이용이 동질적이고 가치형성요인 중 지역요인을 공유하는 지역을 말한다.
\end{choices}
\end{question}
\answer{5}
